\section{Work plan}\label{sec:workplan}

The plan assumes a full-time, 12-month project (\cref{fig:gantt}). Items that need ordering are listed
in \cref{sec:resources} and should be ordered in month~1, because optics, gas fittings and sensors
can have lead times of several weeks. A compressed 6-month version is given at the end of this section.

\begin{figure}[t]
  \centering
  % 12-month work plan in half-month units (1 unit = ~2 weeks).
\begin{ganttchart}[
    x unit=0.43cm, y unit chart=0.5cm, y unit title=0.5cm,
    hgrid={*1{draw=gray!20}}, vgrid={*1{draw=gray!12},*1{draw=gray!35}},
    time slot format=simple, canvas/.append style={draw=none},
    title/.append style={fill=uvviolet!12,draw=gray!40}, title label font=\sffamily\scriptsize\bfseries,
    bar/.append style={fill=uvviolet!55,draw=uvviolet!80!black}, bar height=0.55,
    bar label font=\sffamily\scriptsize, bar label node/.append style={align=right},
    group/.append style={fill=tealdark,draw=tealdark}, group label font=\sffamily\scriptsize\bfseries,
    group right shift=0, group top shift=0.35, group height=0.3, group peaks height=0.15,
    milestone/.append style={fill=amberdark,draw=amberdark!70!black}, milestone height=0.55,
    milestone label font=\sffamily\scriptsize\itshape\color{amberdark!80!black},
    milestone label node/.append style={align=right},
    link/.style={-Stealth,draw=gray!60}
  ]{1}{24}
  \gantttitle{Month}{24} \\
  \gantttitlelist{1,...,12}{2} \\
  \ganttgroup{Set-up}{1}{10} \\
  \ganttbar{WP0 Induction, safety, reading}{1}{3} \\
  \ganttbar{WP1 Forensics of used fibres}{2}{4} \\
  \ganttbar{WP2 Testbed + data logging}{3}{8} \\
  \ganttbar{WP3 End-face cells: design, print, qualify}{4}{10} \\
  \ganttmilestone{MS1 Plan + risk assessments signed off}{2} \\
  \ganttmilestone{MS2 Forensic report; MS3 testbed qualified}{8} \\
  \ganttgroup{Experiments}{7}{22} \\
  \ganttbar{WP4 Baseline in air: power + dose law}{7}{10} \\
  \ganttbar{WP5 Campaign I: atmosphere screening}{11}{14} \\
  \ganttbar{WP6 Campaign II: confirmation + dose-response}{15}{18} \\
  \ganttbar{WP7 Ex-situ surface analysis}{12}{20} \\
  \ganttbar{WP8 Demonstration on the 369.5\,nm system}{17}{22} \\
  \ganttmilestone{MS4 Cells qualified ($<\!100$\,ppm \Oxy{} for 2 weeks)}{10} \\
  \ganttmilestone{MS5 Campaign I decision}{14} \\
  \ganttmilestone{MS6 Campaign II result}{18} \\
  \ganttgroup{Write-up}{8}{24} \\
  \ganttbar[bar/.append style={fill=uvviolet!25}]{WP9 Methods chapter (while building)}{8}{12} \\
  \ganttbar{WP9 Analysis, thesis, paper draft}{18}{24} \\
  \ganttmilestone{MS7 Thesis submitted}{24}
\end{ganttchart}

  \caption{Work packages (WP) and milestones (MS) for a 12-month project in half-month resolution. The
  methods chapter is drafted while the testbed is being built; the analysis and thesis writing overlap
  with the demonstration phase.}
  \label{fig:gantt}
\end{figure}

\begin{table}[t]
\centering\small
\caption{Work packages, main outputs and milestones.}
\label{tab:wps}
\begin{tabularx}{\textwidth}{@{}l Y L{5.0cm}@{}}
\toprule
WP & Content & Output / milestone \\
\midrule
WP0 & inductions (laser, gas, chemical, workshop), literature, project plan & MS1: plan, risk assessments
      and order list agreed (end of month~1) \\
WP1 & forensics of degraded fibres from the lab (\cref{sec:forensics}) & MS2: short forensic report with
      images; choice of extended arms \\
WP2 & testbed: source, splitting, 4--6 arms, detectors, logger, sentinel arm & MS3: testbed qualified
      (drift of sentinel $<\SI{0.3}{\percent}$/day; logging robust for one week) \\
WP3 & end-face cells: CAD, printing, sealing, leak and \Oxy{} hold tests, material screening
    & MS4: cells qualified ($<\SI{100}{ppm}$ \Oxy{} for two weeks, leak rate $<\SI{1e-5}{mbar.L/s}$) \\
WP4 & Phase 1 baseline in air: power and dose dependence, fibre scatter & degradation law; run length
      for later campaigns \\
WP5 & Phase 2 / Campaign I: \arm{AIR}, \arm{CDA}, \arm{N2}, \arm{ABS} & MS5: decision on Campaign~II
      (decision box in \cref{sec:campaigns}) \\
WP6 & Phase 3 / Campaign II: replicates, extended arms, crossed input/output & MS6: confirmatory result
      against pre-registered criteria \\
WP7 & ex-situ surface analysis of selected ends & identification of films or damage \\
WP8 & Phase 4: cells on a working \SI{369.5}{nm} link for $\geq 8$ weeks & demonstration result,
      practical guidelines \\
WP9 & methods chapter (months 4--6), analysis, thesis, paper draft & MS7: thesis submitted \\
\bottomrule
\end{tabularx}
\end{table}

\subsection{Decision points}
\begin{description}[style=nextline,leftmargin=1.2em,itemsep=3pt]
  \item[D1, end of WP4: is the degradation fast enough to study?] If \arm{AIR} at \SI{3}{mW} loses less
        than about \SI{5}{\percent} in two weeks, the campaigns run at \SI{10}{mW}. If necessary the
        intensity is raised further with a smaller mode-field fibre or by running two arms in series on
        one source. If the testbed does not reproduce the lab degradation at all, the difference
        between testbed and lab (wavelength, contaminant sources near the experiments, cleaning
        practice) becomes a research question in itself, and the forensic results guide the next steps.
  \item[D2, end of WP5: which mechanism?] Follow the branches of the decision box in \cref{sec:campaigns}.
  \item[D3, end of WP6: what is deployed?] The cell type (purged or static, with or without oxygen) for
        the Phase~4 demonstration follows from the confirmatory result.
\end{description}

\subsection{Compressed 6-month variant}
If the project has to fit into six months, keep the logic but drop breadth:
month~1: inductions, forensics, ordering; month~2: two to four arms and simple purged cells
(no static cell yet); month~3: short baseline at \SI{10}{mW}; months~4--5: one campaign with \arm{AIR},
\arm{CDA}, \arm{N2} (plus \arm{ABS} if ready); month~6: cleaning sequence, analysis, writing. The
demonstration on a working link and the extended arms become recommendations for a follow-up project.

\begin{keybox}[Minimum viable thesis]
Even if the full testbed is delayed, the following is enough for a good thesis: the forensic analysis
of used fibres, plus one well-controlled parallel comparison of ambient air, clean dry air and dry
\Ntwo{} at the end faces, with the cleaning sequence at the end. The extended arms, the static absorber
cells and the demonstration are the parts that turn it into an excellent thesis and a paper.
\end{keybox}
